\documentclass[11pt]{article}

\usepackage[margin=25mm]{geometry}

\usepackage{amsmath,amssymb,fontspec,unicode-math}

\IfFontExistsTF{DejaVu Serif}{\setmainfont{DejaVu Serif}}{\setmainfont{Latin Modern Roman}}

\IfFontExistsTF{Latin Modern Math}{\setmathfont{Latin Modern Math}}{\setmathfont{DejaVu Math TeX Gyre}}

\usepackage{xurl,hyperref}

\setlength{\parindent}{0pt}\setlength{\parskip}{6pt}

\begin{document}

\title{Imagined Autocurricula}

\author{ELVIA preliminary AI analysis}\date{21 October 2026}\maketitle

Preliminary assessment: Minor.\par

Independent scientific review remains pending. Author adoption has not been requested.

The latest paper still gives two inequivalent replay procedures and two different curriculum priority functions. Exact small witnesses show an actor-gradient sign reversal for stored-action replay and a priority-ranking reversal for accumulated versus plain positive TD errors. These are Minor written-method inconsistencies; the actual training path remains unverified.

\section*{Review scope}

Latest-version, claim-focused review of the cited author-paper passages and applicable current author-release paths. The stated counterexamples and scope controls were checked analytically; no author program, old verifier, training run or benchmark was executed.

\section*{Problem 1: Initial-condition replay and stored-action replay are different actor updates}

Source: Sections 2.4 and 3.3, PDF pp. 4–5; Appendix C, Algorithm 2, PDF p. 19.

Assessment: supported; impact: minor

\subsection*{Source claim (paraphrase)}

The body says current-policy rollouts restart from sampled contexts; Algorithm 2 instead loads a saved trajectory and updates its stored actions without behavior correction.

\subsection*{Analysis}

In a one-step bandit, let rewards be \(1,0\), current good-action probability \(0.2\), replay good-action probability \(0.9\), and action-independent critic baseline \(2\). The true current-policy logit gradient is \(0.2(0.8)=0.16\). Reusing stored actions instead gives \(0.9(0.8)(-1)+0.1(-0.2)(-2)=-0.68\). An action-independent baseline cancels under on-policy sampling; it need not cancel under the replay law. Terminal masking permits this one-step case within a longer maximum imagination horizon.

\subsection*{Finding and scope}

The two written algorithms are not generally equivalent. The witness does not show that the actual experiments used old actions or encountered this state; original implementation/logs are needed for that determination.

\subsection*{Proposed correction}

Make the pseudocode match initial-context replay with fresh current-policy actions, or specify and justify the off-policy alternative. No performance-preserving drop-in repair is established.

The freshly retrieved proceedings and latest target arXiv v2 retain the conflicting body and Algorithm 2 descriptions. The current author project page still says Code (coming soon); no released implementation was available to decide the historical execution path.

\section*{Problem 2: Accumulated positive TD residuals and their mean reverse rankings}

Source: Eq. (7), PDF p. 5; Appendix C, Algorithm 2, PDF p. 19; horizon interpretation, PDF pp. 8 and 20.

Assessment: supported; impact: minor

\subsection*{Source claim (paraphrase)}

Eq.(7) sums future discounted positive TD residuals for every starting position; Algorithm 2 takes their plain mean.

\subsection*{Analysis}

For horizon \(H=5\), \(q=0.985\cdot0.95=0.93575\), and positive TD residual vectors \(A=(0.9,0,0,0,0)\), \(B=(0,0,0.75,0,0)\), plain means are \(0.18,0.15\). The accumulated future-residual score instead gives \(0.18\) and \(0.15(1+q+q^2)=0.421706709375\), reversing their ranking. Binary terminal rewards and value estimates in \([0,1]\) can realize these vectors; the witness does not require illegal reward timing.

\subsection*{Finding and scope}

The priorities are different functions, even at equal horizon; endpoint indexing cannot resolve the reversed ordering. A horizon trend under the accumulated form can partly reflect its position weights, rather than independently demonstrating difficulty.

\subsection*{Proposed correction}

Specify which of the two priority functions is intended. Changing horizon normalization also changes the curriculum; no benchmark-equivalent substitution is asserted.

The freshly retrieved proceedings and latest target arXiv v2 retain the conflicting body and Algorithm 2 descriptions. The current author project page still says Code (coming soon); no released implementation was available to decide the historical execution path.

\section*{Limits}

No world model, policy or old verifier was executed. The current author project page still lists code as forthcoming. The mathematical witnesses distinguish the written procedures but do not identify which procedure generated the reported training results.\par

\section*{Source identity}

NeurIPS 2025 official proceedings PDF (current publisher bytes, checked 2026-10-21)

\url{https://papers.nips.cc/paper_files/paper/2025/file/a6a6891cf1dfc64d664f086cf5976e93-Paper-Conference.pdf}

Reviewed PDF SHA-256: \nolinkurl{a6e4b4f1a6bc0baf66979ae5ec5b27de5b09df71ce11739844821020826ccbab}

This is the same research work as the predecessor, matched by title, author identity and official publisher channel. Its freshly retrieved exact PDF bytes and real retrieval time are recorded. Current target arXiv metadata and the relevant available author-version passages were checked separately; background works were excluded. Original report/source history is unchanged.

\end{document}
